\section{Comparación con la Literatura}
\label{sec:comp}

\subsection{Tesis vs.\ Dazzi et al.\ (2024)}

Ambos trabajan sobre el mismo ground truth (bump B1 y Thacker T1), lo que
permite comparación directa de arquitectura, aunque los \emph{problemas}
difieren: Dazzi resuelve el problema directo mientras que esta tesis
resuelve el inverso.

\begin{table}[H]
\centering
\caption{Comparación arquitectónica: Dazzi et al.\ (2024) vs.\ esta tesis.}
\label{tab:dazzi-comp}
\begin{tabular}{lll}
\toprule
Característica & Dazzi et al.\ (2024) & Esta tesis \\
\midrule
Tarea & Directa (aprender $h, u, \zb$) & Inversa (recuperar $\zb$) \\
Red & MLP único $\to (h, u, \zb)$ & SolNet + BathNet \\
Parámetros & 543{,}603 & 17{,}923 \\
Forma SWE & Aumentada (4 ecs.) & Primitiva (2 ecs.) \\
$\zb$ temporal & $\partial \zb/\partial t = 0$ como loss & Estructural (sin $t$) \\
Encoding & Raw $(x,t)$ & Fourier features \\
Entrenamiento & Adam (30K) & Adam (12K) + L-BFGS (600) \\
\bottomrule
\end{tabular}
\end{table}

La propuesta (A1) usa cerca de \textbf{30$\times$ menos parámetros} que Dazzi
(17\,923 frente a 543\,603), apoyándose en tres elementos de diseño: (1) el
embedding de Fourier, que compensa la red más pequeña; (2) la imposición
estructural de $\partial \zb / \partial t = 0$, que elimina un hiperparámetro;
y (3) el refinamiento con L-BFGS. Si esa reducción de capacidad mantiene una
precisión comparable es lo que cuantifica el estudio de escalamiento de la
Sección~\ref{sec:arch-comp}. \pendiente{comparación de precisión frente a
Dazzi sobre el bump B1.}

\subsection{Tesis vs.\ Angel et al.\ (2024)}

Angel et al.\ establecen el primer benchmark de inversión de batimetría
con datos experimentales reales; su método adjunto-SWE alcanza un NRMSE de
$\approx$14--16\,\% con 2 sensores.

\begin{table}[H]
\centering
\caption{Comparación de enfoque: Angel et al.\ (2024) vs.\ esta tesis.}
\label{tab:angel-comp}
\begin{tabular}{lll}
\toprule
Propiedad & Angel et al.\ (2024) & Esta tesis \\
\midrule
Método & Adjunto + desc.\ gradiente & PINN \\
Modelo & SWE 1D (Dedalus, $M=68$) & SWE 1D (AD sobre redes) \\
Datos & Reales (canal de oleaje) & Reales (Exp.\ 6) + sintéticos \\
Observaciones & 3 sensores puntuales & Grilla / sensores dispersos \\
Bump & Gaussiano, 200\,mm & Parabólico, 200\,mm \\
$kh$ & 1.3 (intermedio) & $\ll 1$ (SWE exacto) \\
\bottomrule
\end{tabular}
\end{table}

\paragraph{Exp.\ 6 --- comparación con Angel.}
El Exp.~6 (Sección~\ref{sec:exp6}) aplica el PINN propuesto a estos mismos
datos reales; el contraste con el método de restricción \textbf{dura} de
Angel (NRMSE $\approx$14--16\,\%) es el núcleo de la comparación. El método
adjunto impone la SWE como restricción dura mediante un solver directo,
mientras que el PINN de penalización \textbf{blanda} admite la equifinalidad
$\eta = h + \zb$ cuando los sensores están fuera del bump.
\pendiente{recuperación cuantitativa del PINN sobre los datos de Angel y el
diagnóstico blando-vs-duro.}

\subsection{Hallazgo sobre formas de ecuación}

El estudio de ablación (Sección~\ref{sec:ablation}) constituye una
\textbf{contribución original}: ningún trabajo previo compara formas
del residuo SWE específicamente para el problema \emph{inverso}.
Tian et al.\ \cite{tian2025} demuestran la ventaja de la forma
primitiva-conservativa para el problema directo; la ablación de esta tesis
determina qué forma es preferible para la inversión. \pendiente{resultado de
la ablación de formas del residuo.}
